\section{Threats to Validity}
\label{sec:threats}

%% Revised 2026-08-30: every residual the response plan (§6) leaves to the extension is named here
%% in the words the rebuttal uses. 2026-09-13 (plan §13): this section is the ONE home of every
%% caveat the results sections used to restate; design facts §IV states are referenced, not
%% repeated. 2026-09-13 readability pass: long sentences split; no threat added or dropped.

\paragraph{Construct.} ``Correct decomposition'' is a judgement, and no metric here assesses
semantic validity with architects. What we recover is a development-time module view rendered in
C4 notation, not a runtime container model (\S\ref{sec:approach}). Both RQ7
references fell below the raw-$\kappa$ flag; label-free measures and the nesting count show the
disagreement is granularity (\S\ref{sec:rq9}). Container-edge F1 is
curation-constrained, not curation-proof. c2c and adjusted a2a are not reported. The recovery
bands and the adequacy thresholds are post-hoc reading labels, not established boundaries. The
build graph is extracted, not validated against an independent evaluated source: CMake's File
API is the configured project's own evaluated model, cross-checked against CMake's graphviz
output, and the MSBuild-C++, automake, and C\#\ L0 parses are static, so a consistently
extracted graph can be consistently incomplete.

\paragraph{Internal.} The committed curations were authored with an LLM assistant and, on three
C\#\ systems and P-1, with the reference in view (Table~\ref{tab:corpus}). Their scores are
oracle-assisted upper bounds (\S\ref{sec:design}; RQ7 bounds the leakage on eShop only). Two C\#\ references transcribe the projects' naming and manifest
conventions, which the refinement extractors also mine, so a rule keyed on the convention scores
near-perfectly by construction; that score is reported only as a labelled reconstruction ceiling.
OpenCV's GT-doc mirrors the CMake module list and abseil's GT-exp the one-library-per-directory
layout, so their floor scores carry the same convention overlap. The LLM baseline's prompt named
the system and used the real identifiers of publicly documented systems, so training-data
contamination is uncontrolled.
The regime claim is confounded by build technology, granularity, style, size, and reference
provenance, all co-varying with language here. The reference-free diagnostic separates the
build-graph property from the language, but it is post hoc, selected on the nine systems it
separates, with one of its two clauses carried by a single system (\S\ref{sec:rq1}): a
candidate diagnostic, not a validated predictor.

\paragraph{External.} Nine OSS  and one proprietary
system are not the industrial population, and the pipeline favours CMake. %
%; the MSBuild-C++ path has no Tier-A system and is never scored against a reference. 
RQ7 ran on two systems of 19 and 15 projects with three curators overall. Curator variance therefore rests on one pair
(\rqNineEshopCuratorMoJoSpread{} MoJoFM points apart), the
effort numbers are two measurements, and we claim nothing about larger C\#\ systems or manual
modelling.
%, whose comparator arm did not run: the governance advantages over hand
%modelling are stated as properties, not measured. 
Blinding removed copying and tuning, not the
priors of shared public documentation (\S\ref{sec:design}). The drift results are largely
construction-verified. The seeded and combined arms confirm correctness under evidence loss, and
the real-history arm covers \rqSixHistPairs{} release pairs of \rqSixHistSystems{} systems with a
handful of findings, under rules that post-date the releases on two of the three. The seeded
operators cover no splits or file moves, and on real history most changes are intended evolution,
so effectiveness on real erosion at scale is not established and the false-finding rate rests on
two raters' labels of \rqSixHistFindings{} findings.

\paragraph{Conclusion.} With $N\le 9$ paired systems we use non-parametric tests with Holm
correction, paired effect sizes, and per-system deltas. We fixed every threshold in the plan
before the data except the adequacy rule and the recovery bands; the plan's dated commit history
is the only record of that order, and no external registry was used. The RQ2 omnibus is reported
with and without the ninth system (\trStats{}).
